% ---------- פרק 2: הגדרת הבעיה, שאלת המחקר והיעדים ----------
% 22.11.2026: גוף שלישי וקיצור. ערכי היעד (90%, שעה, 4 מתוך 5) אושרו ע"י עמית 13.11; 80% ב-REQ-PER-03 והגדרת ספק מוכח – הצעות Claude.
\chapter{הגדרת הבעיה, שאלת המחקר והיעדים}\label{ch:problem}

\section{המצב הרצוי}

הבעיה היא \textbf{איכות ההחלטה ומשך הבדיקה} בשלב טרום ההתקשרות עם ספק חדש בתעשייה הביטחונית-אווירית. טבלה~\ref{tab:goal} מגדירה את המצב הרצוי בחמישה ממדים.

{\small
\begin{longtable}{|>{\raggedright\arraybackslash}p{3.2cm}|>{\raggedright\arraybackslash}p{10.8cm}|}
\caption{המצב הרצוי בסוף הפרויקט}\label{tab:goal}\\
\hline
\textbf{ממד} & \textbf{המצב הרצוי} \\ \hline
\endfirsthead
\hline
\textbf{ממד} & \textbf{המצב הרצוי} \\ \hline
\endhead
משך ההערכה & הערכה מלאה ומנומקת מתוך מידע ציבורי בפחות משעה; מה שאפשר לאמת רק אצל הספק מסומן לבדיקה ידנית \\ \hline
איכות ואמינות & פרופיל חוזקות מאותות ציבוריים: הסמכות, לקוחות וחוזים מוכרים, הכרה פומבית והיסטוריית ריקולים \\ \hline
חברות פרטיות & הערכה גם בלי דוחות כספיים, בעזרת מדדים עקיפים, עם ציון גלוי של כמה מהמידע נמצא \\ \hline
שקיפות & מקור לכל טענה, ורשימה מפורשת של מה שלא נמצא \\ \hline
דגלי סיכון & זיהוי שיטתי של אירועים מכל סוג, ודירוגם לפי ההשפעה על רציפות האספקה \\ \hline
\end{longtable}
}

\textbf{הפער ומחירו.} בארגונים גדולים, קליטת ספק חדש נמשכת שבועות \cite{hicx2021}, וגם בסופה אין מדד שקוף לשאלה כמה מההחלטה נשען על מידע חסר. ספק חדש מוערך לפני שיש איתו היסטוריית עבודה, ולכן קשה לדעת אם מוצריו איכותיים. הפער חמור במיוחד בחברות פרטיות. \en{Incora}, מפיצה פרטית של חלקי תעופה, פשטה רגל ב-2023 אחרי ירידה בהכנסות, נטל חוב כבד ועמידה בזמני האספקה בכ-50\% בלבד מהמשלוחים \cite{incora2023}; \en{ASCO Industries} הותקפה ב-2019 בכופרה, והייצור שלה נעצר \cite{certeu2019}. בשני המקרים נפגעה רציפות האספקה ללקוחות.

\section{שאלת המחקר וההשערות}

הפרויקט משלב \textbf{שאלת מחקר}, שנבדקת דרך השערות שאפשר להפריך, ו\textbf{יעדי תכן} לפי עקרונות \en{SMART}, שמגדירים את אב-הטיפוס.

\begin{quote}
\textbf{שאלת המחקר: באיזו מידה} מתודולוגיה מבוססת סוכני בינה מלאכותית, שנשענת \textbf{על מידע ציבורי בלבד}, מצליחה \textbf{להעריך את האיכות, האמינות והסיכונים} של ספק ולהפיק המלצת התקשרות מנומקת ושקופה \textbf{בשלב טרום ההתקשרות} בתעשייה הביטחונית-אווירית?
\end{quote}

\textbf{שאלות משנה:}
\begin{enumerate}
\item אילו מדדים ציבוריים בכל אחד משבעת הממדים קשורים בספרות לאיכות ספק או לכשל שלו? ארבעה ממדים בונים את \textbf{פרופיל החוזקות} (איכות ואמינות, התאמה טכנית, עמידה ביעדים ומצב פיננסי), ושלושה את \textbf{דגלי הסיכון} (מצב גיאופוליטי, אירועים חריגים ואבטחת סייבר).
\item על חברה פרטית חסרים בדרך כלל דוחות כספיים, יעדים שפורסמו, דיווחי סייבר (חובה שחלה על חברות ציבוריות בארה"ב \cite{sec2023}) ושקיפות בעלות. אילו מדדים עקיפים זמינים לה, ועד כמה הם מחליפים את המידע החסר?
\item איך לשקלל את הממדים כשחלק מהמידע חסר, ואיך החוסר משפיע על ההמלצה ועל רמת הביטחון בה?
\item האם המתודולוגיה מזהה מראש מקרי כשל ציבוריים, וכמה זמן לפני מועד הכשל?
\end{enumerate}

\textbf{הגדרה אופרטיבית של כשל ספק.} כשל נמדד לפי ההשפעה על רציפות האספקה, ולא לפי גודל ההפסד, בסולם של ארבע רמות:
\begin{enumerate}
\item \textbf{חוסם:} הספק ברשימת סנקציות או בקרת יצוא, או שהוא או בעלי השליטה בו פועלים ממדינה עוינת לישראל. בכלי נקבע שמדינה עוינת היא מדינה שישראל מגדירה כמדינת אויב, או מדינה שאין לה יחסים דיפלומטיים עם ישראל. רשימת המדינות תיקבע מול מקור רשמי בשלב בניית השער (\en{T2.1}).
\item \textbf{כשל אספקה:} אירוע שעצר בפועל ייצור או אספקה, כמו פשיטת רגל, אובדן הסמכה, השבתה בגלל מתקפת סייבר, או מלחמה במדינת הספק.
\item \textbf{אות סיכון:} אירוע שמחליש את הספק בלי עדות לעצירת אספקה, כמו הונאה, תביעה, החלפת הנהלה, פריצת סייבר שלא עצרה ייצור או הסלמה מדינית.
\item \textbf{רעש:} אזכור שלילי שאינו מהותי.
\end{enumerate}
ספק שהגיע לרמה 1 או 2 הוא \textbf{מקרה כשל} לצורך התיקוף. בגלל המידע המסווג שתע"א עוסקת בו, אירוע סייבר חמור שלא נסגר מונע "מומלץ" עד לבדיקת אבטחה ידנית, גם אם לא עצר ייצור.

\textbf{הגדרה אופרטיבית של ספק מוכח.} ספק שפועל בענף התעופה עשר שנים לפחות בלי אירוע ברמה 1 או 2. ההגדרה נשענת על ותק בלבד, ולא על הסמכות או פרסים, כי אלה אותות שהכלי עצמו מודד, ובחירה לפיהם הייתה הופכת את התיקוף למעגלי.

\textbf{סיווג נכון.} להערכה חמש תוצאות אפשריות (פרק~\ref{ch:design}): מומלץ, מומלץ בתנאים, לא מומלץ, פסילה בשער ומידע לא מספיק. סיווג נכון של ספק שכשל הוא "לא מומלץ" או "פסילה", ושל ספק מוכח "מומלץ" בלבד; "מומלץ בתנאים" נספר כשגיאה בשתי הקבוצות, כדי שכלי שעונה "בתנאים" על כל ספק לא ייראה מדויק. "מידע לא מספיק" מדווח בנפרד, כשיעור הימנעות. בנוסף מדווח ההפרש בציון $m$ בין הקבוצות. שיעור הסיווג הנכון מחושב על מקרי כשל שאפשר לצפות ועל ספקים מוכחים; שאר המקרים במאגר בודקים רכיב מסוים של הכלי (טבלה~\ref{tab:cases}).

\textbf{ההשערות} מוצגות בטבלה~\ref{tab:hypotheses}. מדגם התיקוף קטן, כעשרה עד עשרים מקרים, ולכן הן נבדקות כמקרי בוחן ומדווחות תיאורית, בלי מבחן מובהקות. קו הבסיס הפיננסי של \en{H1} הוא מדד \en{Z} של \en{Altman} \cite{altman1968}. אפשר לחשב אותו רק לחברה שמפרסמת דוחות כספיים, ולכן לחברה פרטית קו הבסיס אינו מחזיר תשובה, וההשוואה מדווחת בנפרד לחברות ציבוריות ולחברות פרטיות.

{\small
\begin{longtable}{|p{0.7cm}|>{\raggedright\arraybackslash}p{4.4cm}|>{\raggedright\arraybackslash}p{4.6cm}|>{\raggedright\arraybackslash}p{4.1cm}|}
\caption{ההשערות, אופן הבדיקה ומה יפריך כל אחת}\label{tab:hypotheses}\\
\hline
\textbf{מזהה} & \textbf{השערה} & \textbf{איך תיבדק} & \textbf{מה יפריך אותה} \\ \hline
\endfirsthead
\hline
\textbf{מזהה} & \textbf{השערה} & \textbf{איך תיבדק} & \textbf{מה יפריך אותה} \\ \hline
\endhead
\en{H1} & הערכה בשני חלקים מבחינה בין ספקים שכשלו לספקים מוכחים טוב יותר מבדיקה פיננסית בלבד & בדיקה לאחור על שתי הקבוצות מול קו בסיס פיננסי; לכל כשל נמדד גם זמן ההתרעה & שיעור הסיווג הנכון אינו גבוה מקו הבסיס, או שהפרש $m$ אינו חיובי \\ \hline
\en{H2} & בחברות פרטיות הכיסוי נמוך יותר, אבל המדדים העקיפים עדיין מבחינים בין הקבוצות & השוואה בין כשלים פרטיים לספקים מוכחים פרטיים & המדדים העקיפים אינם מבחינים בין הקבוצות \\ \hline
\en{H3} & ההסבר, שני החלקים ורשימת "לבדוק ידנית" שימושיים להחלטה & שאלון בסולם 1--5 על דוחות לדוגמה, למנחה התעשייתי ולאנשי רכש נוספים & ציון ממוצע נמוך מ-4 \\ \hline
\end{longtable}
}

\textbf{יעדי תכן (\en{SMART}).} \textbf{ספציפי:} אב-טיפוס שמפיק דוח הערכה לפי שבעת הממדים, עם שער חסימה, שני חלקים, ציון כיסוי והמלצה מדורגת. \textbf{מדיד:} דוח מלא לכל ספק במדגם התיקוף, עם מקור לכל טענה. \textbf{בר-השגה:} מידע ציבורי ומודלי שפה זמינים. \textbf{רלוונטי:} עונה על הפער שלעיל. \textbf{תחום בזמן:} אב-טיפוס בסוף שבוע 10 ותיקוף בסוף שבוע 13 של תוכנית העבודה (פרק~\ref{ch:workplan}).

\section{מפרט הדרישות}

לכל דרישה מוגדרים מזהה, ערך יעד ושיטת אימות (טבלה~\ref{tab:requirements}).

{\small
\begin{longtable}{|>{\raggedright\arraybackslash}p{2.35cm}|>{\raggedright\arraybackslash}p{4.6cm}|>{\raggedright\arraybackslash}p{1.3cm}|>{\raggedright\arraybackslash}p{2.7cm}|>{\raggedright\arraybackslash}p{2.4cm}|}
\caption{מפרט הדרישות}\label{tab:requirements}\\
\hline
\textbf{מזהה} & \textbf{הדרישה} & \textbf{סיווג} & \textbf{ערך יעד} & \textbf{אימות} \\ \hline
\endfirsthead
\hline
\textbf{מזהה} & \textbf{הדרישה} & \textbf{סיווג} & \textbf{ערך יעד} & \textbf{אימות} \\ \hline
\endhead
\en{REQ-FUN-01} & שער חסימה לפני הציון: סנקציות, בקרת יצוא, מדינת מושב ובעלות & תפקודית & 100\% מהישויות שזוהו; בעלות לא ידועה: "לא נבדק" ובדיקה ידנית & ישויות שידוע שהן ברשימות \\ \hline
\en{REQ-FUN-02} & ציון 0--100 ומשקל לכל ממד; משקל של ממד חסר מחולק מחדש באופן גלוי & תפקודית & 7 מתוך 7 ממדים, עם ציון או "אין מספיק מידע" & בדיקת פלט \\ \hline
\en{REQ-FUN-03} & ציון כיסוי מידע לכל ממד ובסך הכול & תפקודית & בכל דוח & בדיקת פלט \\ \hline
\en{REQ-FUN-04} & סיווג זמינות המידע: ציבורית (\en{A}), פרטית חלקית (\en{B}), כמעט ללא מידע (\en{C}) & תפקודית & 90\% התאמה לסיווג ידני & מדגם התיקוף \\ \hline
\en{REQ-FUN-05} & רשימת נושאים לבדיקה ידנית & תפקודית & בכל דוח עם כיסוי חלקי & בדיקת פלט \\ \hline
\en{REQ-FUN-06} & זיהוי פתוח של אירועים שליליים, גם מסוג שאינו ברשימה & תפקודית & לכל אירוע: קטגוריה, חומרה 1--4, ביטחון ומקור & מקרי התיקוף \\ \hline
\en{REQ-FUN-07} & חומרה לפי השפעה על רציפות האספקה & תפקודית & רמה 1: פסילה; רמה 2 שלא נסגרה: "לא מומלץ" & מקרי התיקוף \\ \hline
\en{REQ-FUN-08} & ממד סייבר: הסמכות, אירועים פומביים, ודירוג חיצוני רק אם פורסם & תפקודית & ציון, כיסוי ורשימה ידנית & מקרי סייבר ציבוריים \\ \hline
\en{REQ-FUN-09} & המלצה בשלוש רמות ושתי תוצאות עצירה (פסילה, מידע לא מספיק), עם הסבר לכל ממד ולכל חלק & תפקודית & 100\% מהדוחות & בדיקת פלט \\ \hline
\en{REQ-FUN-10} & ממד איכות: הסמכות (למשל \en{AS9100}), לקוחות וחוזים, ריקולים והוראות תקינות & תפקודית & ציון, כיסוי ורשימה ידנית & הספקים המוכחים \\ \hline
\en{REQ-PER-01} & הבחנה בין ספקים שכשלו לספקים מוכחים & ביצועים & טובה מקו בסיס פיננסי (\en{H1}) & בדיקה לאחור \\ \hline
\en{REQ-PER-02} & זמן הפקת דוח לספק & ביצועים & פחות משעה & מדידה על מדגם התיקוף \\ \hline
\en{REQ-PER-03} & זיהוי האירועים הידועים במקרי התיקוף & ביצועים & 80\% לפחות & מול רשימה ידנית שנבנית מראש \\ \hline
\en{REQ-NFR-01} & מקור לכל טענה עובדתית & הסבריות & 100\% מהטענות & ביקורת על מדגם \\ \hline
\en{REQ-NFR-02} & ההחלטה הסופית בידי אדם & אילוץ & אין החלטה אוטומטית מלבד השער & סקירת תהליך \\ \hline
\en{REQ-CON-01} & מידע ציבורי בלבד & אילוץ & אפס מקורות פנימיים & ביקורת מקורות \\ \hline
\en{REQ-CON-02} & איסוף פסיבי, בלי סריקת מערכות הספק & אילוץ חוקי & אפס סריקות & סקירת קוד \\ \hline
\en{REQ-CON-03} & אפשרות להריץ בעתיד על מודל מקומי & אילוץ אבטחה & בלי תלות בספק ענן & סקירת ארכיטקטורה \\ \hline
\end{longtable}
}
